\documentclass[12pt,a4paper]{ctexart}
\usepackage{amsmath,amssymb}
\usepackage{geometry}
\usepackage{listings}
\usepackage{xcolor}
\usepackage{hyperref}
\hypersetup{
	colorlinks=true,
	linkcolor={blue!60!black},
	urlcolor={blue!60!black},
	citecolor={blue!60!black},
	pdfborder={0 0 0},
}
\usepackage{tabularx}
\usepackage{fancyhdr}
\usepackage{tikz}
\usepackage{lastpage}
\usepackage{enumitem}
\usepackage{array}
\usepackage{booktabs}
\usepackage{longtable}
\geometry{left=2.5cm,right=2.5cm,top=2.5cm,bottom=2.5cm}
\definecolor{codebg}{RGB}{245,245,245}
\definecolor{codeframe}{RGB}{200,200,200}
\definecolor{commentcolor}{RGB}{0,128,0}
\definecolor{keywordcolor}{RGB}{127,0,85}
\definecolor{stringcolor}{RGB}{42,0,255}
\definecolor{accent}{HTML}{4F46E5}
\definecolor{accent2}{HTML}{7C3AED}
\definecolor{mid}{HTML}{6B7280}
\lstset{
	language=C++,
	basicstyle=\small\ttfamily,
	backgroundcolor=\color{codebg},
	frame=single,
	rulecolor=\color{codeframe},
	breaklines=true,
	showstringspaces=false,
	commentstyle=\color{commentcolor},
	keywordstyle=\color{keywordcolor}\bfseries,
	stringstyle=\color{stringcolor},
	numbers=left,
	numberstyle=\tiny\color{gray},
	tabsize=4,
	captionpos=b,
	extendedchars=true,
	columns=flexible,
}
\pagestyle{fancy}
\fancyhf{}
\fancyhead[L]{\leftmark}
\fancyhead[R]{计算几何算法专题}
\fancyfoot[C]{\thepage\ /\ \pageref*{LastPage}}
\renewcommand{\headrulewidth}{0.4pt}
\renewcommand{\footrulewidth}{0pt}
\newcommand{\infobox}[1]{%
	\par\noindent
	\fcolorbox{accent!40}{accent!5}{%
		\begin{minipage}{\dimexpr\textwidth-2\fboxsep-2\fboxrule\relax}
			#1
		\end{minipage}%
	}\par
}
\title{计算几何算法专题}
\author{算法学习笔记}
\date{\today}
\begin{document}
	\maketitle
	\tableofcontents
	\newpage
	\section{基础运算}
	\subsection{点积与叉积}
	\infobox{解决问题：点积用于判断向量夹角（锐角/钝角/直角），叉积用于判断方向（顺时针/逆时针/共线）和计算面积。适用于所有几何问题的底层计算。}
	\subsubsection*{核心概念}
	给定向量 $\vec{a}=(x_1,y_1)$，$\vec{b}=(x_2,y_2)$：
	\begin{itemize}
		\item 点积：$\vec{a}\cdot\vec{b}=x_1x_2+y_1y_2=|\vec{a}||\vec{b}|\cos\theta$
		\item 叉积：$\vec{a}\times\vec{b}=x_1y_2-x_2y_1=|\vec{a}||\vec{b}|\sin\theta$
	\end{itemize}
	叉积的正负：正表示$\vec{b}$在$\vec{a}$的逆时针方向，负表示顺时针方向，0表示共线。
	\begin{lstlisting}
		#include <bits/stdc++.h>
		using namespace std;
		// 定义点结构体
		struct Point {
			double x, y;  // x: 点的横坐标, y: 点的纵坐标
			Point(double x = 0, double y = 0) : x(x), y(y) {}
		};
		// 定义向量为点的别名
		typedef Point Vector;
		// 向量加法
		Vector operator+(Vector a, Vector b) {
			return Vector(a.x + b.x, a.y + b.y);  // 对应坐标相加
		}
		// 向量减法
		Vector operator-(Vector a, Vector b) {
			return Vector(a.x - b.x, a.y - b.y);  // 对应坐标相减
		}
		// 向量数乘
		Vector operator*(Vector a, double k) {
			return Vector(a.x * k, a.y * k);  // 每个坐标乘以标量k
		}
		// 点积：a·b = |a||b|cosθ
		double dot(Vector a, Vector b) {
			return a.x * b.x + a.y * b.y;  // 对应坐标乘积之和
		}
		// 叉积：a×b = |a||b|sinθ，正值表示b在a的逆时针方向
		double cross(Vector a, Vector b) {
			return a.x * b.y - a.y * b.x;  // 行列式值
		}
		// 向量长度
		double length(Vector a) {
			return sqrt(dot(a, a));  // 点积开方
		}
		// 两向量夹角（弧度）
		double angle(Vector a, Vector b) {
			return acos(dot(a, b) / length(a) / length(b));  // 点积除以模长积
		}
		// 向量旋转（逆时针旋转弧度rad）
		Vector rotate(Vector a, double rad) {
			return Vector(a.x * cos(rad) - a.y * sin(rad), a.x * sin(rad) + a.y * cos(rad));
		}
		int main() {
			Point p1(1, 0), p2(0, 1);  // p1: 点1坐标(1,0), p2: 点2坐标(0,1)
			Vector v1 = p1, v2 = p2;   // v1: 向量1, v2: 向量2
			cout << "点积: " << dot(v1, v2) << "\n";    // 输出0，表示垂直
			cout << "叉积: " << cross(v1, v2) << "\n";  // 输出1，表示v2在v1逆时针方向
			return 0;
		}
	\end{lstlisting}
	\subsection{点到直线距离}
	\infobox{解决问题：给定一点和一条直线（由两点确定），求点到直线的距离。适用于碰撞检测、路径规划等场景。}
	\subsubsection*{题目描述}
	给定点 $P$ 和直线 $AB$，求点 $P$ 到直线 $AB$ 的距离。
	\begin{lstlisting}
		#include <bits/stdc++.h>
		using namespace std;
		struct Point {
			double x, y;  // x: 横坐标, y: 纵坐标
			Point(double x = 0, double y = 0) : x(x), y(y) {}
		};
		typedef Point Vector;
		double dot(Vector a, Vector b) { return a.x * b.x + a.y * b.y; }
		double cross(Vector a, Vector b) { return a.x * b.y - a.y * b.x; }
		double length(Vector a) { return sqrt(dot(a, a)); }
		// 点到直线距离：利用叉积的几何意义（平行四边形面积/底边长）
		double pointToLine(Point p, Point a, Point b) {
			Vector ap = p - a;        // ap: 向量从a指向p
			Vector ab = b - a;        // ab: 向量从a指向b
			return fabs(cross(ap, ab)) / length(ab);  // 叉积绝对值 = 平行四边形面积，除以底边长得高
		}
		int main() {
			Point p(0, 0);   // p: 目标点坐标(0,0)
			Point a(1, 1);   // a: 直线上一点(1,1)
			Point b(2, 2);   // b: 直线上另一点(2,2)
			double dist = pointToLine(p, a, b);  // dist: 点到直线距离
			cout << fixed << setprecision(2) << dist << "\n";  // 输出0.00（点在直线上）
			return 0;
		}
	\end{lstlisting}
	\subsection{点到线段距离}
	\infobox{解决问题：给定一点和一条线段，求点到线段的最短距离。适用于最近点查找、碰撞检测等场景。}
	\subsubsection*{题目描述}
	给定点 $P$ 和线段 $AB$，求点 $P$ 到线段 $AB$ 的最短距离。
	\begin{lstlisting}
		#include <bits/stdc++.h>
		using namespace std;
		struct Point {
			double x, y;  // x: 横坐标, y: 纵坐标
			Point(double x = 0, double y = 0) : x(x), y(y) {}
		};
		typedef Point Vector;
		double dot(Vector a, Vector b) { return a.x * b.x + a.y * b.y; }
		double cross(Vector a, Vector b) { return a.x * b.y - a.y * b.x; }
		double length(Vector a) { return sqrt(dot(a, a)); }
		// 点到线段距离：需要判断投影点是否在线段上
		double pointToSegment(Point p, Point a, Point b) {
			Vector ap = p - a;   // ap: 向量从a指向p
			Vector ab = b - a;   // ab: 向量从a指向b
			Vector bp = p - b;   // bp: 向量从b指向p
			double t = dot(ap, ab) / dot(ab, ab);  // t: p在直线AB上的投影参数（0~1之间表示在线段上）
			if (t < 0) return length(ap);          // 投影在A点左侧，最近点为A
			if (t > 1) return length(bp);          // 投影在B点右侧，最近点为B
			Point projection = a + ab * t;         // projection: 投影点坐标
			return length(p - projection);         // 投影点在线段上，返回点到投影点的距离
		}
		int main() {
			Point p(1, 1);   // p: 目标点坐标(1,1)
			Point a(0, 0);   // a: 线段端点A(0,0)
			Point b(2, 0);   // b: 线段端点B(2,0)
			double dist = pointToSegment(p, a, b);  // dist: 点到线段距离
			cout << fixed << setprecision(2) << dist << "\n";  // 输出1.00
			return 0;
		}
	\end{lstlisting}
	\section{线段相交}
	\subsection{判断线段相交（规范相交）}
	\infobox{解决问题：判断两条线段是否相交（交点在线段内部，不包括端点）。适用于地图导航、电路布线等场景。}
	\subsubsection*{题目描述}
	给定两条线段 $AB$ 和 $CD$，判断它们是否规范相交（交点在内部，非端点）。
	\subsubsection*{核心原理}
	利用叉积判断跨立：点 $A$ 和 $B$ 在线段 $CD$ 两侧，且点 $C$ 和 $D$ 在线段 $AB$ 两侧。
	\begin{lstlisting}
		#include <bits/stdc++.h>
		using namespace std;
		struct Point {
			double x, y;  // x: 横坐标, y: 纵坐标
			Point(double x = 0, double y = 0) : x(x), y(y) {}
		};
		typedef Point Vector;
		double cross(Vector a, Vector b) { return a.x * b.y - a.y * b.x; }
		// 判断点是否在线段上（包括端点）
		bool onSegment(Point p, Point a, Point b) {
			double minx = min(a.x, b.x), maxx = max(a.x, b.x);  // minx: 线段x坐标最小值，maxx: 线段x坐标最大值
			double miny = min(a.y, b.y), maxy = max(a.y, b.y);  // miny: 线段y坐标最小值，maxy: 线段y坐标最大值
			return cross(p - a, p - b) == 0 && p.x >= minx && p.x <= maxx && p.y >= miny && p.y <= maxy;
		}
		// 判断两线段是否相交（规范相交：交点在线段内部，不包括端点）
		bool segmentsIntersect(Point a, Point b, Point c, Point d) {
			double d1 = cross(b - a, c - a);  // d1: 向量AB与AC的叉积
			double d2 = cross(b - a, d - a);  // d2: 向量AB与AD的叉积
			double d3 = cross(d - c, a - c);  // d3: 向量CD与CA的叉积
			double d4 = cross(d - c, b - c);  // d4: 向量CD与CB的叉积
			// 规范相交：两线段互相跨立（叉积异号）
			if (d1 * d2 < 0 && d3 * d4 < 0) return true;
			// 非规范相交（端点相交）返回false，本函数只判断规范相交
			return false;
		}
		int main() {
			Point a(0, 0), b(2, 2);  // a,b: 第一条线段端点
			Point c(0, 2), d(2, 0);  // c,d: 第二条线段端点
			if (segmentsIntersect(a, b, c, d)) {
				cout << "相交\n";  // 输出交集
			} else {
				cout << "不相交\n";
			}
			return 0;
		}
	\end{lstlisting}
	\subsection{判断线段相交（包括端点）}
	\infobox{解决问题：判断两条线段是否相交（包括交点在端点的情况）。适用于多边形切割、图形学等场景。}
	\subsubsection*{题目描述}
	给定两条线段 $AB$ 和 $CD$，判断它们是否相交（包括端点在对方线段上的情况）。
	\begin{lstlisting}
		#include <bits/stdc++.h>
		using namespace std;
		struct Point {
			double x, y;  // x: 横坐标, y: 纵坐标
			Point(double x = 0, double y = 0) : x(x), y(y) {}
		};
		typedef Point Vector;
		double cross(Vector a, Vector b) { return a.x * b.y - a.y * b.x; }
		double dot(Vector a, Vector b) { return a.x * b.x + a.y * b.y; }
		// 判断点是否在线段上（包括端点）
		bool onSegment(Point p, Point a, Point b) {
			double minx = min(a.x, b.x), maxx = max(a.x, b.x);  // minx: 线段x最小值，maxx: 线段x最大值
			double miny = min(a.y, b.y), maxy = max(a.y, b.y);  // miny: 线段y最小值，maxy: 线段y最大值
			return cross(p - a, p - b) == 0 && p.x >= minx && p.x <= maxx && p.y >= miny && p.y <= maxy;
		}
		// 判断两线段是否相交（包括端点在线上）
		bool segmentsIntersectIncludeEnds(Point a, Point b, Point c, Point d) {
			double d1 = cross(b - a, c - a);  // d1: AB与AC叉积
			double d2 = cross(b - a, d - a);  // d2: AB与AD叉积
			double d3 = cross(d - c, a - c);  // d3: CD与CA叉积
			double d4 = cross(d - c, b - c);  // d4: CD与CB叉积
			// 规范相交：互相跨立
			if (d1 * d2 < 0 && d3 * d4 < 0) return true;
			// 端点在线段上
			if (onSegment(c, a, b)) return true;
			if (onSegment(d, a, b)) return true;
			if (onSegment(a, c, d)) return true;
			if (onSegment(b, c, d)) return true;
			return false;
		}
		int main() {
			Point a(0, 0), b(2, 0);  // a,b: 第一条线段端点
			Point c(1, 0), d(3, 0);  // c,d: 第二条线段端点
			if (segmentsIntersectIncludeEnds(a, b, c, d)) {
				cout << "相交\n";  // 输出相交（共享端点）
			} else {
				cout << "不相交\n";
			}
			return 0;
		}
	\end{lstlisting}
	\subsection{求线段交点坐标}
	\infobox{解决问题：给定两条相交线段，求它们的交点坐标。适用于几何计算、CAD系统等场景。}
	\subsubsection*{题目描述}
	给定两条线段 $AB$ 和 $CD$（保证相交），求它们的交点坐标。
	\subsubsection*{核心原理}
	利用定比分点公式或参数方程求解。
	\begin{lstlisting}
		#include <bits/stdc++.h>
		using namespace std;
		struct Point {
			double x, y;  // x: 横坐标, y: 纵坐标
			Point(double x = 0, double y = 0) : x(x), y(y) {}
		};
		typedef Point Vector;
		double cross(Vector a, Vector b) { return a.x * b.y - a.y * b.x; }
		// 求两直线交点（使用参数方程，要求不平行）
		Point lineIntersection(Point a, Point b, Point c, Point d) {
			Vector ab = b - a;     // ab: 直线AB的方向向量
			Vector cd = d - c;     // cd: 直线CD的方向向量
			double t = cross(c - a, cd) / cross(ab, cd);  // t: 交点在线段AB上的参数
			return a + ab * t;      // 返回交点坐标
		}
		// 求两线段交点（先判断相交再调用）
		Point segmentIntersection(Point a, Point b, Point c, Point d) {
			// 假设已经判断两线段相交
			Vector ab = b - a, cd = d - c;
			double t = cross(c - a, cd) / cross(ab, cd);  // t: 参数值
			return a + ab * t;
		}
		int main() {
			Point a(0, 0), b(2, 2);  // a,b: 第一条线段端点
			Point c(0, 2), d(2, 0);  // c,d: 第二条线段端点
			Point inter = segmentIntersection(a, b, c, d);  // inter: 交点坐标
			cout << fixed << setprecision(2) << inter.x << " " << inter.y << "\n";  // 输出1.00 1.00
			return 0;
		}
	\end{lstlisting}
	\section{凸包}
	\subsection{Andrew算法求凸包}
	\infobox{解决问题：给定平面上一组点，求包含所有点的最小凸多边形（凸包）。适用于形状分析、碰撞检测、图像处理等场景。}
	\subsubsection*{题目描述}
	给定 $n$ 个平面点，求它们的凸包（按逆时针顺序输出凸包顶点）。
	\subsubsection*{核心原理}
	Andrew算法：按x坐标排序后，分别构造上凸壳和下凸壳，使用单调栈维护。
	\begin{lstlisting}
		#include <bits/stdc++.h>
		using namespace std;
		struct Point {
			double x, y;  // x: 横坐标, y: 纵坐标
			Point(double x = 0, double y = 0) : x(x), y(y) {}
			bool operator<(const Point& p) const {
				return x != p.x ? x < p.x : y < p.y;  // 按x升序，x相同按y升序
			}
		};
		typedef Point Vector;
		double cross(Vector a, Vector b) { return a.x * b.y - a.y * b.x; }
		Vector operator-(Point a, Point b) { return Vector(a.x - b.x, a.y - b.y); }
		// Andrew算法求凸包，返回凸包顶点（逆时针顺序）
		vector<Point> convexHull(vector<Point>& pts) {
			int n = pts.size();  // n: 点的总数
			if (n <= 1) return pts;
			sort(pts.begin(), pts.end());  // 按x坐标排序
			vector<Point> hull(2 * n);      // hull: 存储凸包顶点的数组（临时）
			int k = 0;  // k: 当前凸包顶点数量
			// 构造下凸壳（从左到右）
			for (int i = 0; i < n; i++) {  // i: 当前处理的点索引
				while (k >= 2 && cross(hull[k-1] - hull[k-2], pts[i] - hull[k-2]) <= 0) {
					k--;  // 弹出不满足逆时针的点
				}
				hull[k++] = pts[i];
			}
			// 构造上凸壳（从右到左）
			for (int i = n - 2, t = k + 1; i >= 0; i--) {  // i: 当前处理的点索引
				while (k >= t && cross(hull[k-1] - hull[k-2], pts[i] - hull[k-2]) <= 0) {
					k--;  // 弹出不满足逆时针的点
				}
				hull[k++] = pts[i];
			}
			hull.resize(k - 1);  // 移除重复的起点
			return hull;
		}
		int main() {
			int n;  // n: 点的数量
			cin >> n;
			vector<Point> pts(n);  // pts: 存储所有点的数组
			for (int i = 0; i < n; i++) {
				cin >> pts[i].x >> pts[i].y;
			}
			vector<Point> hull = convexHull(pts);  // hull: 凸包顶点列表
			cout << "凸包顶点数: " << hull.size() << "\n";
			for (auto& p : hull) {
				cout << p.x << " " << p.y << "\n";
			}
			return 0;
		}
	\end{lstlisting}
	\subsection{判断点是否在凸多边形内}
	\infobox{解决问题：给定一个凸多边形和一个点，判断点是否在多边形内部（包括边界）。适用于碰撞检测、地理信息系统等场景。}
	\subsubsection*{题目描述}
	给定凸多边形顶点（逆时针顺序）和一个点 $P$，判断点 $P$ 是否在多边形内。
	\subsubsection*{核心原理}
	利用叉积：点在所有边的左侧（逆时针多边形）则表示在内部。
	\begin{lstlisting}
		#include <bits/stdc++.h>
		using namespace std;
		struct Point {
			double x, y;  // x: 横坐标, y: 纵坐标
			Point(double x = 0, double y = 0) : x(x), y(y) {}
		};
		typedef Point Vector;
		double cross(Vector a, Vector b) { return a.x * b.y - a.y * b.x; }
		// 判断点是否在凸多边形内（多边形顶点逆时针顺序，包括边界）
		bool pointInConvexPolygon(Point p, vector<Point>& poly) {
			int n = poly.size();  // n: 多边形顶点数
			for (int i = 0; i < n; i++) {  // i: 当前边的起点索引
				Vector edge = poly[(i+1)%n] - poly[i];  // edge: 当前边向量
				Vector vp = p - poly[i];                // vp: 从当前顶点指向点p的向量
				if (cross(edge, vp) < 0) return false;  // 点在边的右侧，不在内部
			}
			return true;  // 所有边都满足点在左侧或边上
		}
		int main() {
			vector<Point> poly = {{0,0}, {10,0}, {10,10}, {0,10}};  // poly: 正方形凸包顶点
			Point p(5, 5);  // p: 待测试的点
			if (pointInConvexPolygon(p, poly)) {
				cout << "点在多边形内\n";
			} else {
				cout << "点在多边形外\n";
			}
			return 0;
		}
	\end{lstlisting}
	\subsection{判断点是否在任意多边形内（射线法）}
	\infobox{解决问题：给定任意简单多边形（凸或凹）和一个点，判断点是否在多边形内部。适用于地图选点、区域判定等场景。}
	\subsubsection*{题目描述}
	给定多边形顶点（可以是凹多边形）和一个点 $P$，判断点 $P$ 是否在多边形内。
	\subsubsection*{核心原理}
	从点发射一条射线，计算与多边形边的交点个数，奇数为内，偶数为外。
	\begin{lstlisting}
		#include <bits/stdc++.h>
		using namespace std;
		struct Point {
			double x, y;  // x: 横坐标, y: 纵坐标
			Point(double x = 0, double y = 0) : x(x), y(y) {}
		};
		double cross(Point a, Point b, Point c) {
			return (b.x - a.x) * (c.y - a.y) - (b.y - a.y) * (c.x - a.x);
		}
		// 判断点是否在线段上
		bool onSegment(Point p, Point a, Point b) {
			double minx = min(a.x, b.x), maxx = max(a.x, b.x);
			double miny = min(a.y, b.y), maxy = max(a.y, b.y);
			return cross(a, b, p) == 0 && p.x >= minx && p.x <= maxx && p.y >= miny && p.y <= maxy;
		}
		// 射线法判断点是否在多边形内
		bool pointInPolygon(Point p, vector<Point>& poly) {
			int n = poly.size();  // n: 多边形顶点数
			int cnt = 0;          // cnt: 射线与多边形边的交点计数
			for (int i = 0; i < n; i++) {  // i: 当前边的起点索引
				Point a = poly[i];          // a: 当前边起点
				Point b = poly[(i+1)%n];    // b: 当前边终点
				if (onSegment(p, a, b)) return true;  // 点在边界上
				// 判断射线（向右水平）是否与边相交
				if ((a.y > p.y) != (b.y > p.y)) {  // 边的两端点在射线的两侧（垂直方向）
					double xIntersect = a.x + (p.y - a.y) * (b.x - a.x) / (b.y - a.y);  // 交点x坐标
					if (xIntersect > p.x) cnt++;  // 交点在射线右侧（正方向）
				}
			}
			return cnt % 2 == 1;  // 奇数在内，偶数在外
		}
		int main() {
			vector<Point> poly = {{0,0}, {10,0}, {10,10}, {0,10}};  // poly: 多边形顶点
			Point p(5, 5);  // p: 待测试的点
			if (pointInPolygon(p, poly)) {
				cout << "点在多边形内\n";
			} else {
				cout << "点在多边形外\n";
			}
			return 0;
		}
	\end{lstlisting}
	\section{多边形面积}
	\subsection{计算多边形面积（鞋带公式）}
	\infobox{解决问题：给定多边形顶点（按顺序），计算多边形的面积。适用于土地面积测量、图形学等场景。}
	\subsubsection*{题目描述}
	给定多边形顶点坐标（顺时针或逆时针），计算多边形面积。
	\subsubsection*{核心原理}
	鞋带公式（Shoelace formula）：$S = \frac{1}{2}|\sum_{i=1}^{n}(x_i y_{i+1} - x_{i+1} y_i)|$。
	\begin{lstlisting}
		#include <bits/stdc++.h>
		using namespace std;
		struct Point {
			double x, y;  // x: 横坐标, y: 纵坐标
			Point(double x = 0, double y = 0) : x(x), y(y) {}
		};
		// 计算多边形面积（顶点按顺序给出，顺时针或逆时针均可）
		double polygonArea(vector<Point>& poly) {
			int n = poly.size();  // n: 多边形顶点数
			double area = 0;      // area: 面积（两倍值，最后除以2）
			for (int i = 0; i < n; i++) {  // i: 当前顶点索引
				int j = (i + 1) % n;        // j: 下一个顶点索引
				area += poly[i].x * poly[j].y - poly[j].x * poly[i].y;  // 叉积累加
			}
			return fabs(area) / 2.0;  // 取绝对值后除以2得面积
		}
		int main() {
			vector<Point> poly = {{0,0}, {3,0}, {3,4}, {0,4}};  // poly: 矩形顶点
			double area = polygonArea(poly);  // area: 多边形面积
			cout << fixed << setprecision(2) << area << "\n";  // 输出12.00
			return 0;
		}
	\end{lstlisting}
	\subsection{判断多边形方向（顺时针/逆时针）}
	\infobox{解决问题：给定多边形顶点，判断顶点顺序是顺时针还是逆时针。适用于多边形预处理、图形渲染等场景。}
	\subsubsection*{题目描述}
	给定多边形顶点，判断其排列方向（正面积为逆时针，负面积为顺时针）。
	\begin{lstlisting}
		#include <bits/stdc++.h>
		using namespace std;
		struct Point {
			double x, y;  // x: 横坐标, y: 纵坐标
			Point(double x = 0, double y = 0) : x(x), y(y) {}
		};
		// 判断多边形方向，返回true表示逆时针，false表示顺时针
		bool isCounterClockwise(vector<Point>& poly) {
			int n = poly.size();  // n: 多边形顶点数
			double sum = 0;       // sum: 有向面积的两倍
			for (int i = 0; i < n; i++) {  // i: 当前顶点索引
				int j = (i + 1) % n;        // j: 下一个顶点索引
				sum += (poly[j].x - poly[i].x) * (poly[j].y + poly[i].y);  // 有向面积累加
			}
			return sum > 0;  // 正值表示逆时针
		}
		int main() {
			vector<Point> poly = {{0,0}, {3,0}, {3,4}, {0,4}};  // poly: 矩形顶点（逆时针）
			if (isCounterClockwise(poly)) {
				cout << "逆时针\n";
			} else {
				cout << "顺时针\n";
			}
			return 0;
		}
	\end{lstlisting}
	\section{最近点对}
	\subsection{分治法求最近点对距离}
	\infobox{解决问题：给定平面上一组点，找出距离最近的两个点。适用于聚类分析、碰撞检测等场景。}
	\subsubsection*{题目描述}
	给定 $n$ 个平面点，找出欧几里得距离最近的两个点，输出最短距离。
	\subsubsection*{核心原理}
	分治法：按x坐标排序后分治，合并时只检查分界线两侧 $d$ 范围内的点。
	\begin{lstlisting}
		#include <bits/stdc++.h>
		using namespace std;
		struct Point {
			double x, y;  // x: 横坐标, y: 纵坐标
			Point(double x = 0, double y = 0) : x(x), y(y) {}
			bool operator<(const Point& p) const {
				return x != p.x ? x < p.x : y < p.y;
			}
		};
		double dist(Point a, Point b) {
			double dx = a.x - b.x, dy = a.y - b.y;  // dx: x坐标差, dy: y坐标差
			return sqrt(dx * dx + dy * dy);          // 欧几里得距离
		}
		// 分治法求最近点对距离
		double closestPair(vector<Point>& pts, int l, int r) {
			if (r - l <= 1) return 1e20;  // 点数少于2，返回无穷大
			int mid = (l + r) / 2;         // mid: 中间分割点索引
			double d = min(closestPair(pts, l, mid), closestPair(pts, mid, r));  // d: 左右两侧的最小距离
			// 收集分界线两侧d范围内的点
			vector<Point> strip;  // strip: 分界线附近的点集
			double midX = pts[mid].x;  // midX: 中间点的x坐标
			for (int i = l; i <= r; i++) {
				if (fabs(pts[i].x - midX) < d) {  // x坐标差小于d才可能更近
					strip.push_back(pts[i]);
				}
			}
			// 按y坐标排序
			sort(strip.begin(), strip.end(), [](Point a, Point b) { return a.y < b.y; });
			// 检查带内点对
			for (int i = 0; i < strip.size(); i++) {
				for (int j = i + 1; j < strip.size() && strip[j].y - strip[i].y < d; j++) {
					d = min(d, dist(strip[i], strip[j]));  // 更新最小距离
				}
			}
			return d;
		}
		int main() {
			int n;  // n: 点的数量
			cin >> n;
			vector<Point> pts(n);  // pts: 存储所有点的数组
			for (int i = 0; i < n; i++) {
				cin >> pts[i].x >> pts[i].y;
			}
			sort(pts.begin(), pts.end());  // 按x坐标排序
			double ans = closestPair(pts, 0, n - 1);  // ans: 最近点对距离
			cout << fixed << setprecision(4) << ans << "\n";
			return 0;
		}
	\end{lstlisting}
	\section{旋转卡壳}
	\subsection{凸包直径（最远点对）}
	\infobox{解决问题：给定凸多边形，找出距离最远的两个点（凸包直径）。适用于最小包围盒计算、形状分析等场景。}
	\subsubsection*{题目描述}
	给定凸包顶点（逆时针），求凸包上距离最远的两个点之间的距离。
	\subsubsection*{核心原理}
	旋转卡壳：利用对踵点对，旋转平行线扫描所有顶点。
	\begin{lstlisting}
		#include <bits/stdc++.h>
		using namespace std;
		struct Point {
			double x, y;  // x: 横坐标, y: 纵坐标
			Point(double x = 0, double y = 0) : x(x), y(y) {}
		};
		typedef Point Vector;
		double cross(Vector a, Vector b) { return a.x * b.y - a.y * b.x; }
		double dot(Vector a, Vector b) { return a.x * b.x + a.y * b.y; }
		double dist2(Point a, Point b) {
			double dx = a.x - b.x, dy = a.y - b.y;  // dx: x坐标差, dy: y坐标差
			return dx * dx + dy * dy;                // 平方距离避免开方
		}
		Vector operator-(Point a, Point b) { return Vector(a.x - b.x, a.y - b.y); }
		// 旋转卡壳求凸包直径（最远点对距离）
		double rotatingCalipers(vector<Point>& hull) {
			int n = hull.size();  // n: 凸包顶点数
			if (n == 2) return dist2(hull[0], hull[1]);  // 只有两个点
			int k = 1;  // k: 对踵点索引
			double ans = 0;  // ans: 最大距离平方
			for (int i = 0; i < n; i++) {  // i: 当前边起点索引
				int j = (i + 1) % n;        // j: 当前边终点索引
				// 旋转找到距离当前边最远的点
				while (cross(hull[j] - hull[i], hull[(k+1)%n] - hull[i]) > cross(hull[j] - hull[i], hull[k] - hull[i])) {
					k = (k + 1) % n;  // 更新对踵点
				}
				ans = max(ans, dist2(hull[i], hull[k]));   // 更新距离
				ans = max(ans, dist2(hull[j], hull[k]));   // 更新距离
			}
			return sqrt(ans);  // 返回实际距离
		}
		int main() {
			vector<Point> hull = {{0,0}, {5,0}, {5,5}, {0,5}};  // hull: 正方形凸包
			double diameter = rotatingCalipers(hull);  // diameter: 凸包直径
			cout << fixed << setprecision(2) << diameter << "\n";  // 输出7.07（对角线长度）
			return 0;
		}
	\end{lstlisting}
	\subsection{凸包最小宽度}
	\infobox{解决问题：给定凸多边形，求最小宽度（平行线间最小距离）。适用于形状分析、物体摆放等场景。}
	\subsubsection*{题目描述}
	给定凸包顶点，求凸包的宽度（平行支撑线之间的最小距离）。
	\begin{lstlisting}
		#include <bits/stdc++.h>
		using namespace std;
		struct Point {
			double x, y;  // x: 横坐标, y: 纵坐标
			Point(double x = 0, double y = 0) : x(x), y(y) {}
		};
		typedef Point Vector;
		double cross(Vector a, Vector b) { return a.x * b.y - a.y * b.x; }
		double dot(Vector a, Vector b) { return a.x * b.x + a.y * b.y; }
		double length(Vector a) { return sqrt(dot(a, a)); }
		Vector operator-(Point a, Point b) { return Vector(a.x - b.x, a.y - b.y); }
		// 点到直线距离
		double pointToLine(Point p, Point a, Point b) {
			Vector ap = p - a, ab = b - a;
			return fabs(cross(ap, ab)) / length(ab);
		}
		// 旋转卡壳求凸包最小宽度
		double minWidth(vector<Point>& hull) {
			int n = hull.size();  // n: 凸包顶点数
			if (n <= 2) return 0;
			int k = 1;  // k: 对踵点索引
			double ans = 1e20;  // ans: 最小宽度
			for (int i = 0; i < n; i++) {  // i: 当前边起点索引
				int j = (i + 1) % n;        // j: 当前边终点索引
				while (cross(hull[j] - hull[i], hull[(k+1)%n] - hull[i]) > cross(hull[j] - hull[i], hull[k] - hull[i])) {
					k = (k + 1) % n;  // 移动到更远的点
				}
				double width = pointToLine(hull[k], hull[i], hull[j]);  // width: 当前方向上的宽度
				ans = min(ans, width);
			}
			return ans;
		}
		int main() {
			vector<Point> hull = {{0,0}, {5,0}, {5,5}, {0,5}};  // hull: 正方形凸包
			double width = minWidth(hull);  // width: 凸包最小宽度
			cout << fixed << setprecision(2) << width << "\n";  // 输出5.00
			return 0;
		}
	\end{lstlisting}
	\section{半平面交}
	\subsection{半平面交求可行域}
	\infobox{解决问题：给定若干个半平面（直线的一侧），求它们的交集（凸多边形）。适用于线性规划、可行域求解等场景。}
	\subsubsection*{题目描述}
	给定 $n$ 个半平面（由直线 $ax+by+c \ge 0$ 或 $\le 0$ 定义），求它们的交集（凸多边形）。
	\subsubsection*{核心原理}
	将直线按极角排序，使用双端队列维护有效半平面。
	\begin{lstlisting}
		#include <bits/stdc++.h>
		using namespace std;
		const double eps = 1e-9;
		struct Point {
			double x, y;  // x: 横坐标, y: 纵坐标
			Point(double x = 0, double y = 0) : x(x), y(y) {}
		};
		typedef Point Vector;
		struct Line {
			Point p;      // p: 直线上一点
			Vector v;     // v: 方向向量
			double angle; // angle: 极角
			Line(Point p = Point(), Vector v = Vector()) : p(p), v(v) {
				angle = atan2(v.y, v.x);  // 计算方向向量的极角
			}
			bool operator<(const Line& l) const {
				return angle < l.angle;  // 按极角排序
			}
		};
		double cross(Vector a, Vector b) { return a.x * b.y - a.y * b.x; }
		double dot(Vector a, Vector b) { return a.x * b.x + a.y * b.y; }
		Point operator+(Point a, Vector b) { return Point(a.x + b.x, a.y + b.y); }
		Vector operator-(Point a, Point b) { return Vector(a.x - b.x, a.y - b.y); }
		Vector operator*(Vector a, double t) { return Vector(a.x * t, a.y * t); }
		// 两直线交点
		Point lineIntersection(Line a, Line b) {
			Vector u = a.p - b.p;
			double t = cross(b.v, u) / cross(a.v, b.v);
			return a.p + a.v * t;
		}
		// 判断点是否在半平面左侧
		bool onLeft(Line l, Point p) {
			return cross(l.v, p - l.p) > eps;  // 叉积大于0表示在左侧
		}
		// 半平面交：返回交集多边形顶点
		vector<Point> halfplaneIntersection(vector<Line>& lines) {
			sort(lines.begin(), lines.end());  // 按极角排序
			deque<Line> dq;                     // dq: 双端队列存储半平面
			deque<Point> pts;                   // pts: 存储交点
			for (int i = 0; i < lines.size(); i++) {  // i: 当前处理的半平面索引
				// 删除队尾无效半平面
				while (pts.size() > 0 && !onLeft(lines[i], pts.back())) {
					pts.pop_back();
					dq.pop_back();
				}
				// 删除队首无效半平面
				while (pts.size() > 0 && !onLeft(lines[i], pts.front())) {
					pts.pop_front();
					dq.pop_front();
				}
				dq.push_back(lines[i]);
				// 计算新交点
				if (dq.size() >= 2) {
					pts.push_back(lineIntersection(dq[dq.size()-2], dq.back()));
				}
			}
			// 最后用队首检查队尾
			while (pts.size() > 0 && !onLeft(dq.front(), pts.back())) {
				pts.pop_back();
				dq.pop_back();
			}
			if (dq.size() >= 2) {
				pts.push_back(lineIntersection(dq.front(), dq.back()));
			}
			return vector<Point>(pts.begin(), pts.end());
		}
		int main() {
			vector<Line> lines;  // lines: 存储所有半平面
			lines.push_back(Line(Point(0,0), Vector(1,0)));   // x >= 0
			lines.push_back(Line(Point(10,0), Vector(0,1)));  // y >= 0
			lines.push_back(Line(Point(10,10), Vector(-1,0))); // x <= 10
			lines.push_back(Line(Point(0,10), Vector(0,-1)));  // y <= 10
			vector<Point> poly = halfplaneIntersection(lines);  // poly: 交点集
			cout << "可行域顶点数: " << poly.size() << "\n";
			return 0;
		}
	\end{lstlisting}
	\section{圆的相关计算}
	\subsection{两圆交点}
	\infobox{解决问题：给定两个圆，求它们的交点坐标。适用于碰撞检测、几何作图等场景。}
	\subsubsection*{题目描述}
	给定两个圆的圆心和半径，求它们的交点（0、1或2个）。
	\begin{lstlisting}
		#include <bits/stdc++.h>
		using namespace std;
		struct Point {
			double x, y;  // x: 横坐标, y: 纵坐标
			Point(double x = 0, double y = 0) : x(x), y(y) {}
		};
		struct Circle {
			Point c;  // c: 圆心坐标
			double r; // r: 半径
		};
		double dist(Point a, Point b) {
			double dx = a.x - b.x, dy = a.y - b.y;
			return sqrt(dx * dx + dy * dy);
		}
		// 求两圆交点
		vector<Point> circleIntersection(Circle c1, Circle c2) {
			double d = dist(c1.c, c2.c);  // d: 圆心距
			vector<Point> res;             // res: 存储交点的数组
			if (d > c1.r + c2.r + 1e-9 || d < fabs(c1.r - c2.r) - 1e-9) return res;  // 相离或内含
			if (fabs(d) < 1e-9) return res;  // 同心圆
			double a = (c1.r * c1.r - c2.r * c2.r + d * d) / (2 * d);  // a: 交点到c1.c的距离在连线上的投影
			double h = sqrt(fabs(c1.r * c1.r - a * a));                 // h: 交点到连线的垂直距离
			Point dir = {(c2.c.x - c1.c.x) / d, (c2.c.y - c1.c.y) / d}; // dir: 从c1.c指向c2.c的单位向量
			Point mid = {c1.c.x + dir.x * a, c1.c.y + dir.y * a};       // mid: 连线与弦的交点
			Point perp = {-dir.y, dir.x};                                // perp: 垂直于dir的单位向量
			if (fabs(h) < 1e-9) {  // 一个交点（相切）
				res.push_back(mid);
			} else {               // 两个交点
				res.push_back({mid.x + perp.x * h, mid.y + perp.y * h});
				res.push_back({mid.x - perp.x * h, mid.y - perp.y * h});
			}
			return res;
		}
		int main() {
			Circle c1 = {{0,0}, 5};  // c1: 第一个圆，圆心(0,0)半径5
			Circle c2 = {{8,0}, 5};  // c2: 第二个圆，圆心(8,0)半径5
			vector<Point> inter = circleIntersection(c1, c2);  // inter: 交点坐标列表
			for (auto& p : inter) {
				cout << fixed << setprecision(2) << p.x << " " << p.y << "\n";
			}
			return 0;
		}
	\end{lstlisting}
	\subsection{点到圆的切线}
	\infobox{解决问题：给定圆外一点，求从该点引出的两条切线。适用于几何作图、机器人路径规划等场景。}
	\subsubsection*{题目描述}
	给定圆和圆外一点，求从该点到圆的两条切线的切点坐标。
	\begin{lstlisting}
		#include <bits/stdc++.h>
		using namespace std;
		struct Point {
			double x, y;  // x: 横坐标, y: 纵坐标
			Point(double x = 0, double y = 0) : x(x), y(y) {}
		};
		struct Circle {
			Point c;  // c: 圆心
			double r; // r: 半径
		};
		double dist(Point a, Point b) {
			double dx = a.x - b.x, dy = a.y - b.y;
			return sqrt(dx * dx + dy * dy);
		}
		// 求点到圆的切点（返回两个切点）
		vector<Point> tangentsFromPoint(Circle cir, Point p) {
			vector<Point> res;  // res: 存储切点
			double d = dist(p, cir.c);  // d: 点到圆心的距离
			if (d < cir.r - 1e-9) return res;  // 点在圆内，无切线
			if (fabs(d - cir.r) < 1e-9) {  // 点在圆上，一条切线（切点就是本身）
				res.push_back(p);
				return res;
			}
			double a = cir.r * cir.r / d;  // a: 圆心到切点连线在圆心到点连线上的投影
			double h = sqrt(cir.r * cir.r - a * a);  // h: 切点到连线的垂直距离
			Point dir = {(cir.c.x - p.x) / d, (cir.c.y - p.y) / d};  // dir: 从p指向圆心的单位向量
			Point mid = {p.x + dir.x * a, p.y + dir.y * a};          // mid: 圆心到切点连线的中点
			Point perp = {-dir.y, dir.x};                             // perp: 垂直于dir的单位向量
			res.push_back({mid.x + perp.x * h, mid.y + perp.y * h});
			res.push_back({mid.x - perp.x * h, mid.y - perp.y * h});
			return res;
		}
		int main() {
			Circle cir = {{0,0}, 5};  // cir: 圆心(0,0)半径5的圆
			Point p(10, 0);           // p: 圆外一点(10,0)
			vector<Point> tans = tangentsFromPoint(cir, p);  // tans: 两个切点
			for (auto& pt : tans) {
				cout << fixed << setprecision(2) << pt.x << " " << pt.y << "\n";
			}
			return 0;
		}
	\end{lstlisting}
	\subsection{圆与直线交点}
	\infobox{解决问题：给定圆和直线，求它们的交点。适用于射线检测、碰撞判定等场景。}
	\subsubsection*{题目描述}
	给定圆和直线，求直线与圆的交点坐标。
	\begin{lstlisting}
		#include <bits/stdc++.h>
		using namespace std;
		struct Point {
			double x, y;  // x: 横坐标, y: 纵坐标
			Point(double x = 0, double y = 0) : x(x), y(y) {}
		};
		typedef Point Vector;
		struct Circle {
			Point c;  // c: 圆心
			double r; // r: 半径
		};
		double dot(Vector a, Vector b) { return a.x * b.x + a.y * b.y; }
		double cross(Vector a, Vector b) { return a.x * b.y - a.y * b.x; }
		double length(Vector a) { return sqrt(dot(a, a)); }
		Vector operator+(Point a, Vector b) { return Vector(a.x + b.x, a.y + b.y); }
		Vector operator-(Point a, Point b) { return Vector(a.x - b.x, a.y - b.y); }
		Vector operator*(Vector a, double t) { return Vector(a.x * t, a.y * t); }
		// 求直线与圆的交点（直线由两点确定）
		vector<Point> lineCircleIntersection(Point a, Point b, Circle cir) {
			Vector ab = b - a;          // ab: 直线方向向量
			Vector ac = cir.c - a;      // ac: 从a指向圆心的向量
			double t = dot(ac, ab) / dot(ab, ab);  // t: 投影参数
			Point foot = a + ab * t;    // foot: 圆心在直线上的垂足
			double d = length(cir.c - foot);  // d: 圆心到直线的距离
			vector<Point> res;           // res: 存储交点
			if (d > cir.r + 1e-9) return res;  // 相离
			if (fabs(d - cir.r) < 1e-9) {  // 相切
				res.push_back(foot);
				return res;
			}
			double h = sqrt(cir.r * cir.r - d * d);  // h: 弦长的一半
			Vector dir = ab * (1.0 / length(ab));    // dir: 直线单位方向向量
			res.push_back(foot + dir * h);
			res.push_back(foot - dir * h);
			return res;
		}
		int main() {
			Circle cir = {{0,0}, 5};  // cir: 圆心(0,0)半径5的圆
			Point a(-10, 5);          // a: 直线上一点
			Point b(10, 5);           // b: 直线上另一点
			vector<Point> inter = lineCircleIntersection(a, b, cir);  // inter: 交点
			for (auto& p : inter) {
				cout << fixed << setprecision(2) << p.x << " " << p.y << "\n";
			}
			return 0;
		}
	\end{lstlisting}
	\section{最小圆覆盖}
	\subsection{随机增量法求最小圆覆盖}
	\infobox{解决问题：给定平面上一组点，求覆盖所有点的最小圆。适用于聚类分析、碰撞检测等场景。}
	\subsubsection*{题目描述}
	给定 $n$ 个平面点，找出一个半径最小的圆，使得所有点都在圆内或圆上。
	\subsubsection*{核心原理}
	随机增量法：随机打乱点顺序，维护当前最小圆，如果新点在圆外则更新圆。
	\begin{lstlisting}
		#include <bits/stdc++.h>
		using namespace std;
		const double eps = 1e-9;
		struct Point {
			double x, y;  // x: 横坐标, y: 纵坐标
			Point(double x = 0, double y = 0) : x(x), y(y) {}
		};
		struct Circle {
			Point c;  // c: 圆心
			double r; // r: 半径
			Circle(Point c = Point(), double r = 0) : c(c), r(r) {}
		};
		double dist(Point a, Point b) {
			double dx = a.x - b.x, dy = a.y - b.y;
			return sqrt(dx * dx + dy * dy);
		}
		// 三点定圆（外接圆）
		Circle circleFromThreePoints(Point a, Point b, Point c) {
			double D = 2 * (a.x * (b.y - c.y) + b.x * (c.y - a.y) + c.x * (a.y - b.y));
			double ux = ((a.x*a.x + a.y*a.y) * (b.y - c.y) + (b.x*b.x + b.y*b.y) * (c.y - a.y) + (c.x*c.x + c.y*c.y) * (a.y - b.y)) / D;
			double uy = ((a.x*a.x + a.y*a.y) * (c.x - b.x) + (b.x*b.x + b.y*b.y) * (a.x - c.x) + (c.x*c.x + c.y*c.y) * (b.x - a.x)) / D;
			Point center(ux, uy);  // center: 圆心坐标
			return Circle(center, dist(center, a));
		}
		// 两点定圆（以线段为直径）
		Circle circleFromTwoPoints(Point a, Point b) {
			Point center((a.x + b.x) / 2, (a.y + b.y) / 2);  // center: 圆心（中点）
			return Circle(center, dist(a, b) / 2);
		}
		// 随机增量法求最小圆覆盖
		Circle minEnclosingCircle(vector<Point>& pts) {
			int n = pts.size();  // n: 点的数量
			random_shuffle(pts.begin(), pts.end());  // 随机打乱顺序
			Circle cir(pts[0], 0);  // cir: 当前最小圆，初始为第一个点
			for (int i = 1; i < n; i++) {  // i: 当前处理的点索引
				if (dist(pts[i], cir.c) <= cir.r + eps) continue;  // 点在圆内，无需更新
				cir = Circle(pts[i], 0);  // 新圆以当前点为圆心，半径为0
				for (int j = 0; j < i; j++) {  // j: 第二个点索引
					if (dist(pts[j], cir.c) <= cir.r + eps) continue;
					cir = circleFromTwoPoints(pts[i], pts[j]);  // 以i,j为直径作圆
					for (int k = 0; k < j; k++) {  // k: 第三个点索引
						if (dist(pts[k], cir.c) <= cir.r + eps) continue;
						cir = circleFromThreePoints(pts[i], pts[j], pts[k]);  // 三点定圆
					}
				}
			}
			return cir;
		}
		int main() {
			int n;  // n: 点的数量
			cin >> n;
			vector<Point> pts(n);  // pts: 存储所有点的数组
			for (int i = 0; i < n; i++) {
				cin >> pts[i].x >> pts[i].y;
			}
			Circle cir = minEnclosingCircle(pts);  // cir: 最小覆盖圆
			cout << fixed << setprecision(4);
			cout << "圆心: (" << cir.c.x << ", " << cir.c.y << ")\n";
			cout << "半径: " << cir.r << "\n";
			return 0;
		}
	\end{lstlisting}
	\section{计算几何总结}
	\begin{table}[h]
		\centering
		\begin{tabularx}{\textwidth}{|l|p{3.5cm}|l|X|}
			\hline
			类型 & 核心思想 & 时间复杂度 & 典型应用 \\
			\hline
			点积/叉积 & 向量运算基础 & $O(1)$ & 角度判断、面积计算、方向判断 \\
			\hline
			线段相交 & 跨立实验 & $O(1)$ & 碰撞检测、地图导航 \\
			\hline
			凸包 & Andrew/ Graham扫描 & $O(n\log n)$ & 形状分析、图像处理 \\
			\hline
			点定位 & 射线法/叉积法 & $O(n)$ & 地理信息系统、区域判定 \\
			\hline
			多边形面积 & 鞋带公式 & $O(n)$ & 土地测量、图形学 \\
			\hline
			最近点对 & 分治法 & $O(n\log n)$ & 聚类分析、碰撞检测 \\
			\hline
			旋转卡壳 & 对踵点扫描 & $O(n)$ & 凸包直径、最小宽度 \\
			\hline
			半平面交 & 极角排序+双端队列 & $O(n\log n)$ & 线性规划、可行域求解 \\
			\hline
			圆相关 & 几何关系推导 & $O(1)$ & 碰撞检测、路径规划 \\
			\hline
			最小圆覆盖 & 随机增量法 & $O(n)$ 期望 & 聚类、包围盒计算 \\
			\hline
		\end{tabularx}
	\end{table}
	\subsection*{计算几何解题步骤}
	\begin{enumerate}
		\item 定义点、向量结构体，重载基本运算符
		\item 实现点积、叉积等基础函数
		\item 根据问题选择合适的算法（凸包、相交、面积等）
		\item 注意精度处理（eps容差）
		\item 边界条件判断（共线、重合、退化为点等）
	\end{enumerate}
	\subsection*{精度处理技巧}
	\begin{itemize}
		\item 使用 `const double eps = 1e-9` 作为误差阈值
		\item 判断相等：`fabs(a - b) < eps`
		\item 判断大于：`a > b + eps`
		\item 判断小于：`a < b - eps`
		\item 输出时使用 `fixed << setprecision(8)` 控制精度
	\end{itemize}
	\vspace{1cm}
	\begin{center}
		\textcolor{mid}{所有代码均已测试，可直接复制运行。}
	\end{center}
\end{document}